%--- iliad ---
% summary: >-
%   What must a next-token predictor remember about the past? Causal states,
%   belief states and their geometry, from HMMs to generalised HMMs, ending
%   with the evidence that transformers learn this structure.
%--- end ---
\documentclass[11pt]{article}
\IfFileExists{iliad.sty}{\usepackage[boxes]{iliad}}{\usepackage[boxes]{../iliad}}
% float: the Mess3 solution carries a figure inside its box, where a float
% cannot go; [H] sets it in place.
\usepackage{float}

\title{Computational Mechanics}
\author{\authorname{Xavier Poncini}\\ \affiliation{Simplex}
  \and \authorname{Adam Shai}\\ \affiliation{Simplex}
  \and \authorname{Paul Riechers}\\ \affiliation{Simplex}}
\date{}

\begin{document}
\maketitle

\section{Prerequisites}\label{prerequisites}

Familiarity with the goals and basic methodology of mechanistic interpretability, 
including the notions of features and circuits; the linear representation hypothesis; 
linear probes as a method for reading off internal representations; and a basic 
understanding of the transformer architecture. 

Students should be comfortable with the idea that one can train a linear map 
from model activations to some target structure and evaluate its 
quality (e.g. via MSE or \(R^2\)).

\textbf{Mathematical background.} The module is mathematically self-contained: 
all formal definitions (HMMs, GHMMs, belief states, MSPs) are introduced from scratch. 
However, students will engage with the material more fluently if they are comfortable with the following:

\begin{itemize}
\item
  \textbf{Linear algebra:} row-stochastic matrices, rank and invertibility, null spaces (kernels).
\item
  \textbf{Probability:} conditional probability, Bayes\textquotesingle{} rule, probability distributions over finite sets, the probability simplex.
\item
  \textbf{Basic machine learning:} next-token prediction, loss functions (cross-entropy), the concept of model activations.
\end{itemize}

\begin{learningoutcomes}
\begin{itemize}
\item
  Construct their own hidden Markov models.
\item
  Explain the motivation for introducing generalised hidden Markov models (GHMMs).
\item
  Explain why the belief state is a useful object for making predictions about GHMM data.
\item
  Compute the belief state of a GHMM.
\item
  Explain the difference between the belief state of a HMM and the belief state of a GHMM (that is itself not an HMM).
\item
  Explain the relationship between belief states and next-token observation probabilities for GHMMs.
\item
  Compute the mixed state presentation (MSP) for GHMMs.
\item
  Explain why generating data is easier than predicting data.
\item
  Explain the evidence for the claim: transformers represent belief state geometry in their residual stream.
\item
  Develop hypotheses for where transformers represent: belief states, observation probabilities, log observation probabilities.
\end{itemize}
\end{learningoutcomes}

\section{Content}\label{content}

\subsection{Fast-track}\label{fast-track}

Read \href{https://arxiv.org/pdf/2405.15943}{Transformers represent belief state geometry in their residual stream}. Focus on being able to answer the following questions:

\begin{itemize}
\item
  What is a hidden Markov model (HMM)?
\item
  What is a belief state?
\item
  What is the mixed state presentation (MSP)?
\item
  About the map learnt from activations to belief states:
  \begin{itemize}
  \item
    What is the source?
  \item
    What is the target?
  \item
    How is the quality of the map evaluated?
  \end{itemize}
\end{itemize}

% The roadmap is the Schedule from the source repo's README
% (iliad-team/iliad-intensive-C.3), which is what the day was taught from on
% 10 Sep 2026. Its links pointed at the repo's own PDFs; they point at this
% page's decks and sections instead. Words are the README's.
\subsection{Schedule}\label{schedule}

\begin{itemize}
\item
  \textbf{10:00--10:30} --- \href{https://iliad-intensive.org/downloads/computational-mechanics/computational-mechanics-slides.pdf}{Lecture: Overview and scope}
\item
  \textbf{10:30--11:15} --- \href{https://iliad-intensive.org/downloads/computational-mechanics/computational-mechanics-slides-2-predicting-the-future.pdf}{Lecture: Predicting the future}
\item
  \textbf{11:15--11:30} --- \emph{Break}
\item
  \textbf{11:30--12:30} --- \hyperref[sec:predicting-the-future]{{Exercises}: Predicting the future}
\item
  \textbf{12:30--13:30} --- \emph{Lunch}
\item
  \textbf{13:30--14:15} --- \href{https://iliad-intensive.org/downloads/computational-mechanics/computational-mechanics-slides-3-representing-the-past.pdf}{Lecture: Representing the past}
\item
  \textbf{14:15--15:45} --- \hyperref[sec:representing-the-past]{{Exercises}: Representing the past}
\item
  \textbf{15:45--16:00} --- \emph{Break}
\item
  \textbf{16:00--17:45} --- \hyperref[sec:belief-geometry]{Reading \& Discussion: Transformers represent belief geometry}
\item
  \textbf{17:45--18:00} --- Feedback
\end{itemize}

% Parts 2-4 are the source repo's written-exercises/CompMechExercisesPt{2,3}
% and reading-guides/CompMechReadingGuidePt4, run through a scripted
% transform (environment names only; every word and formula is the
% author's). Part 1 is lecture-only: its deck is the first Slides row.
\section{Predicting the future}\label{sec:predicting-the-future}
% Shared source for the student sheet and the solutions sheet.
% Content inside a solution environment appears only in the solutions PDF.

\begin{callout}[note][Instructions]
    Give a clear justification for each answer. You may use results proved in
    the lecture unless an exercise asks you to establish them. Questions marked
    \textbf{Extension} are optional.
\end{callout}


\begin{exercise}[From statistics to partitions]
\label{ex:p2-1}
  Let $u_1,\ldots,u_6\in\mathcal X_+^*$ be six possible histories. Suppose
  that they induce the following conditional distributions over the entire
  future:
  \[
    \begin{array}{c|cccccc}
      \text{history} & u_1 & u_2 & u_3 & u_4 & u_5 & u_6\\
      \hline
      \text{future distribution} & P & P & P & Q & Q & S
    \end{array}
  \]
  where $P$, $Q$, and $S$ are pairwise distinct.

  Three statistics $R_A$, $R_B$, and $R_C$ induce the partitions
  \[
    \begin{aligned}
      \Pi_A
        &=\bigl\{\{u_1,u_2\},\{u_3\},\{u_4\},\{u_5\},\{u_6\}\bigr\},\\
      \Pi_B
        &=\bigl\{\{u_1,u_2,u_3\},\{u_4,u_5\},\{u_6\}\bigr\},\\
      \Pi_C
        &=\bigl\{\{u_1,u_2,u_3,u_4\},\{u_5\},\{u_6\}\bigr\}.
    \end{aligned}
  \]
  Recall that a statistic is sufficient for prediction when histories assigned
  the same value induce the same distribution over the future.

  \begin{enumerate}
    \item What does it mean for two histories to belong to the same cell of the
      partition induced by a statistic? Illustrate your answer using $u_1$ and
      $u_2$ in $\Pi_A$.

    \item Determine which of $R_A$, $R_B$, and $R_C$ are sufficient. If a
      statistic is not sufficient, give two histories witnessing the failure.

    \item Which cells of $\Pi_A$ can be merged without losing sufficiency?
      Perform all such merges and identify the resulting partition.

    \item Explain why no two cells of $\Pi_B$ can be merged while preserving
      sufficiency. What does this show about $R_B$?

    \item Prove that every refinement of a sufficient partition is also
      sufficient. Is every coarsening of a sufficient partition necessarily
      sufficient? Use $\Pi_A$, $\Pi_B$, and $\Pi_C$ to illustrate your answer.

    \item \textbf{Extension.} Consider the partition
      \[
        \Pi_D
          =\bigl\{\{u_1\},\{u_2,u_3\},\{u_4,u_5\},\{u_6\}\bigr\}.
      \]
      Show that $\Pi_A$ and $\Pi_D$ are both sufficient but that neither
      refines the other. How can both nevertheless refine the same minimal
      sufficient partition?
  \end{enumerate}
\end{exercise}

\begin{exercise}[Causal states are minimal sufficient statistics]
\label{ex:p2-2}
  For $x,x'\in\mathcal X_+^*$, define the relation $\sim$ by
  \[
    \begin{aligned}
    x\sim x'
    \quad\Longleftrightarrow\quad
    &\Pr\!\left(
      X_{|x|+1:|x|+|y|}=y\mid X_{1:|x|}=x
    \right)\\
    &\qquad=
    \Pr\!\left(
      X_{|x'|+1:|x'|+|y|}=y\mid X_{1:|x'|}=x'
    \right)
    \end{aligned}
  \]
  for every $y\in\mathcal X^*$. Define the causal state of $x$ by
  \[
    \epsilon(x):=[x]_{\sim},
  \]
  and write $\mathcal S:=\epsilon(\mathcal X_+^*)$ for the set of
  causal states.

  \begin{enumerate}
    \item Prove that $\sim$ is an equivalence relation on
      $\mathcal X_+^*$ by establishing each of the following:
      \begin{enumerate}
        \item reflexivity: $x\sim x$ for every $x\in\mathcal X_+^*$;
        \item symmetry: if $x\sim x'$, then $x'\sim x$;
        \item transitivity: if $x\sim x'$ and $x'\sim x''$, then
          $x\sim x''$.
      \end{enumerate}

    \item Prove that $\epsilon$ is sufficient for prediction.

    \item Let $R:\mathcal X_+^*\to\mathcal R$ be any sufficient statistic.
      Prove that
      \[
        R(x)=R(x')\quad\Longrightarrow\quad
        \epsilon(x)=\epsilon(x').
      \]
      Deduce that every sufficient partition refines the causal-state
      partition.

    \item Using parts (b) and (c), conclude that $\epsilon$ is a minimal
      sufficient statistic.

    \item \textbf{Extension.} Let $T$ be another minimal sufficient
      statistic. Prove that $T$ and $\epsilon$ induce the same partition of
      $\mathcal X_+^*$. Deduce that there is a bijection
      \[
        \psi:\mathcal S\to\operatorname{im}(T)
      \]
      satisfying $T=\psi\circ\epsilon$.
  \end{enumerate}
\end{exercise}

\begin{exercise}[From finite generators to causal states]
\label{ex:p2-3}
  In this exercise, $S_t$ denotes the hidden state immediately before the
  transition that emits $X_{t+1}$. All parameters lie strictly between $0$
  and $1$.

  \begin{enumerate}
    \item \textbf{The Zero--One--Random Process.}
      Let $S_0$ be uniformly distributed over
      $\{s_0,s_1,s_R\}$. The process has transitions
      \[
        s_0\xrightarrow{\,0:1\,}s_1,
        \qquad
        s_1\xrightarrow{\,1:1\,}s_R,
        \qquad
        s_R
        \begin{cases}
          \xrightarrow{\,0:1-p\,}s_0,\\[-1mm]
          \xrightarrow{\,1:p\,}s_0.
        \end{cases}
      \]
      An edge label $x:q$ means that the process emits $x$ and follows the
      edge with probability $q$.

      \begin{enumerate}
        \item List all possible histories of length two. For each history,
          determine whether it uniquely identifies the current hidden phase.
          When it does, identify that phase.

        \item For each ambiguous length-two history, determine which
          additional observations resolve the ambiguity. Identify the
          resulting phase and explain why it remains known after every
          subsequent observation.

        \item Assume that $\epsilon(\emptyset)$, $\epsilon(0)$,
          $\epsilon(1)$, and $\epsilon(10)$ are pairwise distinct. Show that
          the causal states have shortest-length representatives
          \[
            \emptyset,\quad 0,\quad 1,\quad 10,\quad 00,\quad 01,\quad 11.
          \]
          Explain why no additional causal states arise.

        \item Under the assumption in part (iii), draw the automaton
          associated with the causal-state dynamics.
      \end{enumerate}

    \item \textbf{The Even Process.}
      Let $S_0$ be uniformly distributed over $\{A,B\}$ and consider the
      generator
      \[
        A\xrightarrow{\,0:q\,}A,
        \qquad
        A\xrightarrow{\,1:1-q\,}B,
        \qquad
        B\xrightarrow{\,1:1\,}A.
      \]
      Thus every run of $1$s between successive $0$s has even length.

      \begin{enumerate}
        \item For each of the length-one histories $0$ and $1$, determine
          whether the current hidden state can be identified exactly.

        \item Starting from each ambiguous history in part (i), consider its
          possible one-symbol extensions. Which extensions resolve the
          ambiguity, and which remain ambiguous? You may use the fact that
          $\epsilon(\emptyset)=\epsilon(11)$. What does this imply for
          histories of the form $1^n$?

        \item Assume additionally that
          $\epsilon(\emptyset)\neq\epsilon(1)$. Show that the causal states
          have shortest representatives
          \[
            \emptyset,\qquad 1,\qquad 0,\qquad 01,
          \]
          and explain why there are no additional causal states. You may
          distinguish causal states using possible and impossible future
          continuations.

        \item Using $\epsilon(\emptyset)=\epsilon(11)$ and assuming
          $\epsilon(\emptyset)\neq\epsilon(1)$, draw the automaton associated
          with the causal-state dynamics.
      \end{enumerate}

    \item \textbf{The Simple Nonunifilar Source.}
      Let $S_0$ be uniformly distributed over $\{A,B\}$ and consider
      \[
        A\xrightarrow{\,0:1/2\,}A,
        \qquad
        A\xrightarrow{\,0:1/2\,}B,
        \qquad
        B\xrightarrow{\,0:1/2\,}B,
        \qquad
        B\xrightarrow{\,1:1/2\,}A.
      \]

      \begin{enumerate}
        \item Show that observing a $1$ identifies the new hidden state as
          $A$.

        \item Starting from $A$, show that there are $n+1$ hidden paths that
          emit $0^n$, each with probability $2^{-n}$. How many of these paths
          can subsequently emit $1$? Hence calculate
          \[
            \Pr(X_{n+2}=1\mid X_{1:n+1}=10^n).
          \]

        \item Show that this probability is different for every $n\geq0$.
          Conclude that the histories
          \[
            1,10,100,\ldots
          \]
          belong to distinct causal states and that the process has infinitely
          many causal states.
      \end{enumerate}
  \end{enumerate}
\end{exercise}

\begin{exercise}[Why next-token prediction is enough]
\label{ex:p2-4}
  Recall that
  \[
    \mathcal X_+^*
      :=\left\{x\in\mathcal X^*:
        \Pr\!\left(X_{1:|x|}=x\right)>0\right\}
  \]
  is the set of allowable histories.

  Consider a model with parameters $\theta$ satisfying
  \[
    h_n=f_\theta(h_{n-1},x_n),
    \qquad
    y_n=g_\theta(h_{n-1},x_n).
  \]
  Fix the initial hidden state $h_0$ and suppose that the model predicts the
  next-token distribution exactly: for every $n\geq1$,
  \[
    y_n
      =\Pr\!\left(
        X_{n+1}\mid X_{1:n}=x_{1:n}
      \right),
      \qquad \forall\,x_{1:n}\in\mathcal X_+^*.
  \]
  For every nonempty allowable history, define
  \[
    R(x_{1:n}):=(h_{n-1},x_n).
  \]

  For $a\in\mathcal X$, define the deterministic update
  \[
    \Phi_a(h,x):=\bigl(f_\theta(h,x),a\bigr).
  \]
  For a word $z_{1:j}$, write
  \[
    \Phi_{z_{1:j}}
      :=\Phi_{z_j}\circ\cdots\circ\Phi_{z_1},
    \qquad
    \Phi_{z_{1:0}}:=\operatorname{id}.
  \]

  \begin{enumerate}
    \item Show that, whenever $x_{1:n}z_{1:j}$ is allowable,
      \[
        R(x_{1:n}z_{1:j})
          =\Phi_{z_{1:j}}\!\left(R(x_{1:n})\right).
      \]
      Hence show that, if $R(x_{1:n})=R(x'_{1:m})$ and both extended
      histories are allowable, then
      \[
        R(x_{1:n}z_{1:j})=R(x'_{1:m}z_{1:j}).
      \]

    \item Let $R(x_{1:n})=R(x'_{1:m})$ and fix
      $z_{1:k}\in\mathcal X^k$. As the symbols of $z_{1:k}$ are read from
      left to right, use part (a) and exact next-token prediction to show that
      for each $0\leq j<k$ such that the prefixes through $z_{1:j}$ are
      allowable,
      \[
        \begin{aligned}
        &\Pr\!\left(
          X_{n+j+1}=z_{j+1}
          \mid X_{1:n+j}=x_{1:n}z_{1:j}
        \right)\\
        &\qquad=
        \Pr\!\left(
          X_{m+j+1}=z_{j+1}
          \mid X_{1:m+j}=x'_{1:m}z_{1:j}
        \right).
        \end{aligned}
      \]
      Deduce that either the same first zero-probability extension is
      encountered after both histories, or every prefix of $z_{1:k}$ is
      allowable after both histories.

    \item In the case where every prefix is allowable, use the conditional
      chain rule to write
      \[
        \begin{aligned}
        &\Pr\!\left(
          X_{n+1:n+k}=z_{1:k}\mid X_{1:n}=x_{1:n}
        \right)\\
        &\qquad=
        \prod_{j=0}^{k-1}
        \Pr\!\left(
          X_{n+j+1}=z_{j+1}
          \mid X_{1:n+j}=x_{1:n}z_{1:j}
        \right).
        \end{aligned}
      \]
      Write the analogous product after $x'_{1:m}$ and compare its factors.
      Deal separately with the zero-probability case from part (b), and
      conclude directly that
      \[
        \begin{gathered}
          R(x_{1:n})=R(x'_{1:m})
          \quad\Longrightarrow\\[1mm]
          \Pr\!\left(
            X_{n+1:n+k}=z_{1:k}\mid X_{1:n}=x_{1:n}
          \right)
          =
          \Pr\!\left(
            X_{m+1:m+k}=z_{1:k}\mid X_{1:m}=x'_{1:m}
          \right).
        \end{gathered}
      \]
      Hence $R(x_{1:n})=(h_{n-1},x_n)$ is a sufficient statistic for
      predicting the entire future.
  \end{enumerate}
\end{exercise}


\section{Representing the past}\label{sec:representing-the-past}
% Shared source for the student sheet and the solutions sheet.
% Content inside a solution environment appears only in the solutions PDF.

\begin{callout}[note][Instructions]
    Give a clear justification for each answer. You may use results proved in
    the lecture unless an exercise asks you to establish them. Questions marked
    \textbf{Extension} are optional.
\end{callout}


\begin{exercise}[Translating between HMM diagrams and matrices]
\label{ex:p3-1}
  For an edge-emitting hidden Markov model, use the convention
  \[
    T^{(a)}_{ij}
      :=\Pr(X_{t+1}=a,S_{t+1}=j\mid S_t=i).
  \]
  Thus rows index the current hidden state, columns index the next hidden
  state, and the superscript records the emitted symbol. Throughout this
  exercise, $0<p,q<1$.

  \begin{enumerate}
    \item Suppose that $T^{(a)}_{ij}=r$. Explain each component of the
      following diagram:
      \[
        i\xrightarrow{\,a:r\,}j.
      \]
      What does a zero entry mean? Why is the matrix
      $T:=\sum_{a\in\mathcal X}T^{(a)}$ row-stochastic?

    \item \textbf{From matrices to a diagram: Zero--One--Random.}
      Take the state order to be $(0,1,R)$ and suppose that
      \[
        \eta^{(\emptyset)}
          =\begin{bmatrix}\tfrac13&\tfrac13&\tfrac13\end{bmatrix},
        \qquad
        T^{(0)}=
        \begin{bmatrix}
          0&1&0\\
          0&0&0\\
          1-p&0&0
        \end{bmatrix},
        \qquad
        T^{(1)}=
        \begin{bmatrix}
          0&0&0\\
          0&0&1\\
          p&0&0
        \end{bmatrix}.
      \]
      \begin{enumerate}
        \item Draw the corresponding edge-emitting HMM. Describe informally
          the three phases represented by $0$, $1$, and $R$.

        \item Calculate $\Pr(X_{1:3}=011)$ and
          $\Pr(X_{1:3}=000)$ in two ways: first by tracing paths through your
          diagram, and then from
          \[
            \Pr(X_{1:3}=x_{1:3})
              =\eta^{(\emptyset)}
                T^{(x_1)}T^{(x_2)}T^{(x_3)}\mathbf 1.
          \]
      \end{enumerate}

    \item \textbf{From a diagram to matrices: the Even Process.}
      Consider the edge-emitting HMM
      \[
        A\xrightarrow{\,0:q\,}A,
        \qquad
        A\xrightarrow{\,1:1-q\,}B,
        \qquad
        B\xrightarrow{\,1:1\,}A,
      \]
      with no other edges. Take the state order to be $(A,B)$ and let
      $\eta^{(\emptyset)}=\begin{bmatrix}\tfrac12&\tfrac12\end{bmatrix}$.
      \begin{enumerate}
        \item Construct $T^{(0)}$ and $T^{(1)}$. Account explicitly for every
          zero entry in the two matrices.

        \item Form $T=T^{(0)}+T^{(1)}$ and verify that it is row-stochastic.

        \item Calculate $\Pr(X_{1:3}=011)$ and
          $\Pr(X_{1:3}=010)$ both by tracing paths through the diagram and by
          multiplying the appropriate symbol matrices. Explain how the second
          calculation reflects the defining constraint of the Even Process.
      \end{enumerate}

    \item \textbf{Extension.} Let $\mathcal X$ be the observation alphabet of
      an arbitrary finite-state HMM, let
      $T=\sum_{a\in\mathcal X}T^{(a)}$, and let
      $\eta^{(\emptyset)}$ be a probability distribution. Prove that, for
      every $n\geq 0$,
      \[
        \sum_{w\in\mathcal X^n}\Pr(X_{1:n}=w)
          =\eta^{(\emptyset)}T^n\mathbf 1
          =1.
      \]
  \end{enumerate}
\end{exercise}

\begin{exercise}[Belief states and a lossy next-token readout]
\label{ex:p3-2}
  Continue with the Zero--One--Random HMM from \cref{ex:p3-1}, using state order
  $(0,1,R)$. For an allowable history $x_{1:n}=x_1\cdots x_n$, compute its
  belief from scratch as
  \[
    \eta^{(x_{1:n})}
      :=\frac{\eta^{(\emptyset)}T^{(x_1)}\cdots T^{(x_n)}}
        {\eta^{(\emptyset)}T^{(x_1)}\cdots T^{(x_n)}\mathbf 1}.
  \]
  Here $\eta_i^{(x_{1:n})}=\Pr(S_n=i\mid X_{1:n}=x_{1:n})$. For an allowable
  extension $x_{1:n+1}=x_{1:n}a$, update recursively as
  \[
    \eta^{(x_{1:n+1})}
      =\frac{\eta^{(x_{1:n})}T^{(a)}}
        {\eta^{(x_{1:n})}T^{(a)}\mathbf 1}.
  \]

  \begin{enumerate}
    \item Compute $\eta^{(0)}$ and $\eta^{(1)}$. Interpret their components
      as posterior probabilities over the three hidden phases.

    \item Show that the seven beliefs
      \[
        \mathcal B=
        \bigl\{
          \eta^{(\emptyset)},\eta^{(0)},\eta^{(1)},\eta^{(00)},
          \eta^{(01)},\eta^{(10)},\eta^{(11)}
        \bigr\}
      \]
      are distinct and that $\mathcal B$ is closed under every allowable
      symbol update. Draw the symbol-labelled transition diagram on
      $\mathcal B$.

    \item Compare your belief-state diagram with both the causal-state
      automaton constructed in Part 2 and the three-state generative HMM from
      \cref{ex:p3-1}.
      \begin{enumerate}
        \item Match each belief to the causal state with the same
          shortest-length representative. Do the symbol-labelled transitions
          agree?

        \item Which beliefs correspond to knowing the hidden phase exactly,
          and which represent uncertainty over several phases? Show that the
          hidden phase becomes known after at most three observations and
          remains known thereafter.

        \item Observe that the finite-history belief-state presentation has
          seven states while the generative HMM has only three. What does this
          suggest about the possible relative difficulty of generation and
          inference?
      \end{enumerate}

    \item For each belief $\eta^{(x_{1:n})}\in\mathcal B$ from part (b),
      compute its associated next-token probability vector
      \[
        \begin{gathered}
        p^{(x_{1:n})}
        :=\left(
          \Pr(X_{n+1}=a\mid X_{1:n}=x_{1:n})
        \right)_{a\in\{0,1\}}
        =\eta^{(x_{1:n})}A,\\[1mm]
        A_{ia}:=\sum_jT^{(a)}_{ij}.
        \end{gathered}
      \]
      Visualise the belief geometry on the $2$-simplex over hidden states
      $(0,1,R)$ and the next-token geometry on the $1$-simplex over symbols
      $(0,1)$.

    \item Calculate
      \[
        \Pr(X_{3:4}=00\mid X_{1:2}=10)
        \qquad\text{and}\qquad
        \Pr(X_{3:4}=00\mid X_{1:2}=01).
      \]
      Deduce that the next-token probability vector need not be sufficient for
      predicting the entire future.

    \item \textbf{Extension.} Let
      $\delta:=\eta^{(10)}-\eta^{(01)}$. Show that
      \[
        \delta A=0
        \qquad\text{but}\qquad
        \delta T^{(0)}T^{(0)}\mathbf 1\neq0.
      \]
      Interpret these two statements in terms of one-step and longer-horizon
      prediction.
  \end{enumerate}
\end{exercise}

\begin{exercise}[Nonunifilar generators and large belief-state automata]
\label{ex:p3-3}
  For a belief $\eta$ and a symbol $a$ satisfying
  $\eta T^{(a)}\mathbf 1>0$, write
  \[
    F_a(\eta):=\frac{\eta T^{(a)}}{\eta T^{(a)}\mathbf 1}.
  \]
  Thus the belief-state automaton contains the transition
  \[
    \eta\xrightarrow{\,a:\eta T^{(a)}\mathbf 1\,}F_a(\eta).
  \]

  \begin{enumerate}
    \item \textbf{The Simple Nonunifilar Source.}
      Revisit the source from Part 2. With state order $(A,B)$, its symbol
      matrices and initial belief are
      \[
        T^{(0)}=
        \begin{bmatrix}\tfrac12&\tfrac12\\0&\tfrac12\end{bmatrix},
        \qquad
        T^{(1)}=
        \begin{bmatrix}0&0\\\tfrac12&0\end{bmatrix},
        \qquad
        \eta^{(\emptyset)}=\begin{bmatrix}\tfrac12&\tfrac12\end{bmatrix}.
      \]
      \begin{enumerate}
        \item For $n\geq0$, show that
          \[
            \eta^{(10^n)}=
              \begin{bmatrix}\dfrac1{n+1}&\dfrac n{n+1}\end{bmatrix},
            \qquad
            \eta^{(0^n)}=
              \begin{bmatrix}\dfrac1{n+2}&\dfrac{n+1}{n+2}\end{bmatrix}.
          \]
          Deduce that
          $\eta^{(0^n)}=\eta^{(10^{n+1})}$ and observe that
          $\eta^{(\emptyset)}=\eta^{(10)}$.

        \item Show that, for every $n\geq1$,
          \[
            \eta^{(10^n1)}=\eta^{(1)}.
          \]

        \item Draw the first four states of the belief-state automaton and
          indicate its continuing structure with an ellipsis. Show that the
          beliefs $\eta^{(10^n)}$ are all distinct, identify their limiting
          point in the $1$-simplex, and explain why the two-state generator
          produces a countably infinite belief-state automaton.
      \end{enumerate}

    \item \textbf{Mess3.}
      Now take hidden-state and symbol sets both equal to $\{0,1,2\}$, with
      \[
        \alpha=\frac15,\qquad b=\frac{1-\alpha}{2}=\frac25,
        \qquad x=\frac3{20},\qquad y=1-2x=\frac7{10},
      \]
      and
      \[
        \begin{aligned}
          T^{(0)}&=
          \begin{bmatrix}
            \alpha y&bx&bx\\ \alpha x&by&bx\\ \alpha x&bx&by
          \end{bmatrix},&
          T^{(1)}&=
          \begin{bmatrix}
            by&\alpha x&bx\\ bx&\alpha y&bx\\ bx&\alpha x&by
          \end{bmatrix},\\[1mm]
          T^{(2)}&=
          \begin{bmatrix}
            by&bx&\alpha x\\ bx&by&\alpha x\\ bx&bx&\alpha y
          \end{bmatrix},&
          \eta^{(\emptyset)}&=
          \begin{bmatrix}\tfrac13&\tfrac13&\tfrac13\end{bmatrix}.
        \end{aligned}
      \]
      \begin{enumerate}
        \item Verify that $T=\sum_{a=0}^2T^{(a)}$ is row-stochastic and
          explain why this generator is nonunifilar. Compute
          $\eta^{(0)},\eta^{(1)}$, and $\eta^{(2)}$.

        \item Compute $\eta^{(00)},\eta^{(01)}$, and $\eta^{(02)}$. Use the
          cyclic symmetry of the process to sketch the first two generations
          of the belief-state automaton, including the symbol labels.

        \item Compute the next-token matrix $A$, with
          $A_{ia}=\sum_jT^{(a)}_{ij}$, and show that it is invertible. What
          does this imply about whether distinct beliefs can be merged by the
          map $p^{(x_{1:n})}=\eta^{(x_{1:n})}A$? Contrast this with \cref{ex:p3-2}.

        \item Using a short program, form
          \[
            \mathcal B_d:=
            \bigl\{\eta^{(x_{1:d})}:x_{1:d}\in\{0,1,2\}^d\bigr\},
            \qquad 0\leq d\leq7.
          \]
          Record the number of distinct beliefs at each depth. Visualise
          $\mathcal B_7$ on the hidden-state $2$-simplex, colouring each point
          by its final symbol. Also visualise the associated next-token
          vectors on the symbol $2$-simplex. Describe the geometric structure
          that emerges and how the three maps $F_0,F_1,F_2$ generate the
          outgoing transitions of the belief-state automaton.

        \item Explain why the belief-state automata constructed in this
          exercise are unifilar even though their generators are not. Compare
          the ladder-with-resets structure of the Simple Nonunifilar Source
          with the branching structure of Mess3. Finally, explain why the set
          of beliefs reached by finite histories is countable even though its
          closure can have a much richer fractal geometry.
      \end{enumerate}
  \end{enumerate}
\end{exercise}

\begin{exercise}[Generalised hidden Markov models and predictive vectors]
\label{ex:p3-4}
  A $d$-dimensional generalised hidden Markov model (GHMM) consists of a
  finite alphabet $\mathcal X$, real $d\times d$ matrices
  $(T^{(a)})_{a\in\mathcal X}$, an initial row vector
  $\eta^{(\emptyset)}$, and a column vector $\phi$. Writing
  \[
    T:=\sum_{a\in\mathcal X}T^{(a)},
  \]
  these objects satisfy
  \[
    T\phi=\phi,\qquad
    \eta^{(\emptyset)}\phi=1,\qquad
    \eta^{(\emptyset)}T^{(x_1)}\cdots T^{(x_n)}\phi\geq0
  \]
  for every $x_{1:n}\in\mathcal X^n$ and every $n\geq0$.

  \begin{enumerate}
    \item \textbf{A probability law from matrix products.}
      Consider
      \[
        \Pr(X_{1:n}=x_{1:n})
          =\eta^{(\emptyset)}
            T^{(x_1)}\cdots T^{(x_n)}\phi.
      \]
      \begin{enumerate}
        \item Show that these quantities are nonnegative and that
          \[
            \sum_{x_{1:n}\in\mathcal X^n}
              \Pr(X_{1:n}=x_{1:n})=1
          \]
          for every $n\geq0$.

        \item Show that
          \[
            \sum_{a\in\mathcal X}
              \Pr(X_{1:n+1}=x_{1:n}a)
              =\Pr(X_{1:n}=x_{1:n}).
          \]
          Explain why this is the consistency condition needed to define a
          one-sided stochastic process.

        \item The GHMM conditions do not require
          $\eta^{(\emptyset)}T=\eta^{(\emptyset)}$. Explain why the process
          need not therefore be stationary. Show that this additional
          equality is sufficient for stationarity.
      \end{enumerate}

    \item \textbf{Predictive vectors.}
      For any $x_{1:n}$ satisfying
      $\Pr(X_{1:n}=x_{1:n})>0$, define its predictive vector by
      \[
        \eta^{(x_{1:n})}
          :=\frac{\eta^{(\emptyset)}
              T^{(x_1)}\cdots T^{(x_n)}}
            {\eta^{(\emptyset)}
              T^{(x_1)}\cdots T^{(x_n)}\phi}.
      \]
      \begin{enumerate}
        \item Show that
          \[
            \eta^{(x_{1:n})}\phi=1
          \]
          and, for every continuation $y_{1:m}$,
          \[
            \Pr(X_{n+1:n+m}=y_{1:m}\mid X_{1:n}=x_{1:n})
              =\eta^{(x_{1:n})}
                T^{(y_1)}\cdots T^{(y_m)}\phi.
          \]

        \item Derive the recursive update and identify its normalising
          factor:
          \[
            \eta^{(x_{1:n}a)}
              =\frac{\eta^{(x_{1:n})}T^{(a)}}
                {\eta^{(x_{1:n})}T^{(a)}\phi},
            \quad
            \Pr(X_{n+1}=a\mid X_{1:n}=x_{1:n})
              =\eta^{(x_{1:n})}T^{(a)}\phi.
          \]

        \item Conceptually contrast an HMM belief state with a GHMM
          predictive vector. Compare the meaning of their coordinates, their
          positivity and normalisation properties, and how they encode
          predictions of future observations.
      \end{enumerate}

    \item \textbf{Predictive vectors need not be probability vectors.}
      Consider the two-state HMM
      \[
        M^{(0)}=
          \begin{bmatrix}\tfrac45&0\\0&\tfrac15\end{bmatrix},
        \qquad
        M^{(1)}=
          \begin{bmatrix}\tfrac15&0\\0&\tfrac45\end{bmatrix},
        \qquad
        \alpha=\begin{bmatrix}\tfrac12&\tfrac12\end{bmatrix}.
      \]
      Thus a hidden state is selected initially and thereafter emits a biased
      coin without changing state. Let
      \[
        S=\begin{bmatrix}2&-1\\0&1\end{bmatrix},
        \qquad
        T^{(a)}:=S^{-1}M^{(a)}S,\qquad
        \eta^{(\emptyset)}:=\alpha S,\qquad
        \phi:=S^{-1}\mathbf 1.
      \]
      \begin{enumerate}
        \item Compute $T^{(0)},T^{(1)},\eta^{(\emptyset)}$, and $\phi$.
          Show that, for every $x_{1:n}$,
          \[
            \eta^{(\emptyset)}
              T^{(x_1)}\cdots T^{(x_n)}\phi
              =
            \alpha M^{(x_1)}\cdots M^{(x_n)}\mathbf 1.
          \]
          Deduce that this change of coordinates gives a valid GHMM for the
          same process.

        \item Compute $\eta^{(0)}$ and $\eta^{(1)}$. Verify directly from
          $\eta^{(0)}$ that
          \[
            \Pr(X_2=0\mid X_1=0)=\frac{17}{25},
            \qquad
            \Pr(X_2=1\mid X_1=0)=\frac8{25}.
          \]
          Explain why the entries of a GHMM predictive vector need not
          themselves be probabilities.
      \end{enumerate}

    \item \textbf{The probability matrix and minimal dimension.}
      For finite sets $\mathcal U,\mathcal V\subseteq\mathcal X^*$, define
      the probability matrix by
      \[
        \mathsf P_{\mathcal U,\mathcal V}(u,v)
          :=\Pr(X_{1:|u|+|v|}=uv),
        \qquad u\in\mathcal U,\quad v\in\mathcal V.
      \]
      \begin{enumerate}
        \item Use the GHMM matrix product to factor
          $\mathsf P_{\mathcal U,\mathcal V}$ as a matrix of prefix row
          vectors times a matrix of continuation column vectors. Deduce that
          \[
            \operatorname{rank}\mathsf P_{\mathcal U,\mathcal V}\leq d
          \]
          for every $d$-dimensional GHMM realisation.

        \item For the process in part (c), take
          $\mathcal U=\mathcal V=\{\emptyset,0,1\}$. Compute
          $\mathsf P_{\mathcal U,\mathcal V}$ and use an SVD to find its
          singular values. Deduce that the two-dimensional GHMM is
          minimal-dimensional.
      \end{enumerate}

  \end{enumerate}
\end{exercise}


\section{Identifying belief geometry in transformers}\label{sec:belief-geometry}

So far, we have developed the conceptual framework of computational mechanics.
We now ask whether neural networks trained for next-token prediction---including
transformers and recurrent architectures---learn the predictive structures
identified by that framework. You will read selected sections of two papers,
then discuss their evidence and limitations in small groups.

\begin{callout}[note][Session at a glance]
  \begin{enumerate}
    \item \textbf{Readings (75 minutes).} Begin with the main reading and use
      the route below. Continue to the extension if time permits.
    \item \textbf{Discussion (30 minutes).} Form groups of three or four.
      Compare what you found convincing, confusing, or incomplete; use the
      prompts on page 2 if useful.
  \end{enumerate}
\end{callout}

\subsection{Readings}
Prioritise the main reading. If you can answer its evaluation question, you
have probably understood one of its central results. If time is short, skim
the extension overview before deciding whether to continue.

\subsubsection{[Main] \href{https://arxiv.org/pdf/2405.15943}{Transformers represent belief state geometry in their residual stream}}
  \textbf{Reading route}
\begin{itemize}
    \item \textbf{Skim.} Sections 1, 2.1, and 2.2.
    \item \textbf{Read.} Sections 2.3 and 3--5.
    \item \textbf{Evaluation.} Explain Figure 6D. What do the comparison and
      shuffle control establish?
  \end{itemize}

  \textbf{Overview.} This paper presents evidence that transformers trained on
  HMM-generated sequences encode belief-state geometry in their residual
  streams. The claim is tested on the Mess3 and Random--Random--XOR (RRXOR)
  processes. For each process, the authors fit an affine map from residual
  activations to the exact HMM belief state and evaluate it using mean-squared
  error (MSE). The error decreases over training checkpoints. A shuffle
  control tests whether low MSE could arise merely because a high-dimensional
  activation space is being projected into a low-dimensional belief space.

\subsubsection{[Extension] \href{https://arxiv.org/pdf/2507.07432}{Neural networks leverage nominally quantum and post-quantum representations}}
  \textbf{Reading route}
\begin{itemize}
    \item \textbf{Skim.} Sections 1--3.
    \item \textbf{Read.} Sections 4--8.
    \item \textbf{Evaluation.} Explain Figure 3B. How are the light-blue and light-orange
      plots related to the corresponding dark-blue and dark-orange plots?
  \end{itemize}

  \textbf{Overview.} This paper extends the question from HMMs to generalised
  hidden Markov models (GHMMs). As in Part 3, a GHMM retains the matrix-product
  expression for output probabilities while relaxing the entrywise
  probabilistic constraints on its internal matrices. This larger model class
  can realise some processes in finite dimension even when a corresponding
  minimal HMM requires exponentially many more hidden states.

  Using techniques related to those in the main reading, the authors study
  transformers, LSTMs, RNNs, and GRUs trained on processes with distinct GHMM
  and HMM realisations. Their results indicate that the networks learn
  belief-state representations associated with the lower-dimensional GHMM
  realisation. The finding suggests that neural networks can select among
  representational model classes in ways that capture dimensionality savings.

\subsection{Discussion}
Here are some prompts to consider:
\begin{itemize}
  \item What is the most convincing evidence that transformers represent
    belief-state geometry? What is the least convincing?
  \item What additional experiment would most strengthen or weaken the claim?
  \item Which alternative hypotheses could explain the reported results?
  \item Can you explain each element of the diagram below---including the
    generative HMM, Bayesian updating, transformer activations, fitted affine
    map, and belief-state geometry?
\end{itemize}

\begin{figure}[ht]
  \centering
  \includegraphics[width=0.88\linewidth]{fig/shai-belief-state-geometry.png}
  \label{fig:shai-belief-state-geometry}
\end{figure}


% From the source repo's README, "Python exercises: HMM from scratch".
\section{Python exercises: HMM from scratch}\label{sec:python-exercises}

The Python exercises are ARENA-style notebooks in two parts. Each part has an
\textbf{exercises} notebook (with \texttt{\# YOUR CODE HERE} stubs, inline tests, and
collapsible solutions) and a fully worked \textbf{solutions} notebook.

\begin{itemize}
\item
  \textbf{Part 1 --- Sequence probabilities \& next-token distributions} (exercises) --- \href{https://colab.research.google.com/github/iliad-team/iliad-intensive-C.3/blob/main/python-exercises/colab/part1_sequence_probabilities_exercises_colab.ipynb}{Open in Colab}
\item
  Part 1 (solutions) --- \href{https://colab.research.google.com/github/iliad-team/iliad-intensive-C.3/blob/main/python-exercises/colab/part1_sequence_probabilities_solutions_colab.ipynb}{Open in Colab}
\item
  \textbf{Part 2 --- Belief states} (exercises) --- \href{https://colab.research.google.com/github/iliad-team/iliad-intensive-C.3/blob/main/python-exercises/colab/part2_belief_states_exercises_colab.ipynb}{Open in Colab}
\item
  Part 2 (solutions) --- \href{https://colab.research.google.com/github/iliad-team/iliad-intensive-C.3/blob/main/python-exercises/colab/part2_belief_states_solutions_colab.ipynb}{Open in Colab}
\end{itemize}

The Colab notebooks are self-contained: their setup cells write the helper
modules into the session, so there is nothing else to install or download.

\section{Learn more}\label{learn-more}

\begin{itemize}
\item
  Beyond toy architectures -- in-context learning
  \begin{itemize}
  \item
    LLMs have an amazing ability to adapt internal representations in-context -- {[}Read: ICLR: \href{https://arxiv.org/pdf/2501.00070}{In-Context Learning of Representations}{]}
  \item
    LLMs can predict HMM data in-context -- {[}Read: \href{https://arxiv.org/pdf/2506.07298}{Pre-trained Large Language Models Learn Hidden Markov Models In-context}{]}
  \item
    This suggests that internal representations of LLMs predicting GHMM data in-context, may resemble belief states.
  \end{itemize}
\item
  Processes consisting of GHMMs with richer structure {[}Read: \href{https://arxiv.org/pdf/2602.02385v1}{Transformers learn factored representations}{]}
\item
  Review article of computational mechanics {[}Read: \href{https://www.nature.com/articles/nphys2190}{Between order and chaos}{]}
\item
  \href{https://www.simplexaisafety.com/}{Simplex} is a non-profit research organisation working on applying computational mechanics to AI safety.
\end{itemize}

\end{document}
